\chapter{Ngại Hỏi Khi Không Hiểu}
\chapterintrobi{Không hiểu nhưng sợ hỏi, rồi mang cả chỗ rối về nhà}{Be afraid to ask}{Đây là thói quen khiến khó khăn nhỏ kịp lớn lên thành nút thắt thật sự.}

\begin{lucbatbox}{Lục bát: Cứ để lát nữa}
Đoạn kia em thấy chưa thông\
Miệng thì muốn hỏi mà lòng lại ngăn\
Sợ rằng lộ vẻ lăn tăn\
Sợ người bên cạnh nghĩ rằng mình chậm\
Tới khi về tối mới thấm\
Một câu không hỏi thành trăm mối mù
\end{lucbatbox}

\begin{tutuyetbox}{Thất ngôn tứ tuyệt: Im lặng vì ngại}
Không hiểu nhưng em vẫn gật đầu\
Sợ hỏi rồi ai cũng biết mau\
Nếu dũng cảm hỏi ngay khi chưa rõ\
Ngày mai đâu phải rối thêm sâu
\end{tutuyetbox}

\begin{ghichu}
Ngại hỏi là tâm lý dễ hiểu. Nhưng trong học tập, nhiều khi một câu hỏi sớm cứu được rất nhiều giờ bối rối muộn.
\end{ghichu}

\begin{englishnote}
\textbf{ask questions}: đặt câu hỏi.\par
\textbf{confused}: chưa hiểu, rối.\par
\textbf{I was too shy to ask.}: Em ngại nên không dám hỏi.
\end{englishnote}